%% ─────────────────────────────────────────────────────────────────────────────
%% CV — Talvin Ackbaraly (English Version) — NVIDIA, Developer Technology intern
%% Serif (XCharter), one restrained accent colour, thin rules, bullets, no icons
%% No photo (requested by NVIDIA): the header runs full width.
%% Compile: latexmk -pdf cv-talvin-ackbaraly-en.tex
%% ─────────────────────────────────────────────────────────────────────────────
\documentclass[10.5pt, a4paper]{extarticle}

\usepackage[T1]{fontenc}
\usepackage[utf8]{inputenc}
\usepackage[english]{babel}
\usepackage[top=1cm, bottom=1cm, left=1.6cm, right=1.6cm]{geometry}
\usepackage{XCharter}
\usepackage[letterspace=80]{microtype}
\usepackage{titlesec}
\usepackage{enumitem}
\usepackage{tabularx}
\usepackage{ragged2e}
\usepackage{xcolor}
\usepackage[scaled=0.94,noupquote]{inconsolata}
\usepackage{tikz}
\usepackage{varwidth}
\usepackage[hidelinks]{hyperref}

\hypersetup{pdftitle={Resume — Talvin Ackbaraly}, pdfauthor={Talvin Ackbaraly}}

%% ── Colour system: one accent, everything else neutral ───────────────────────
%% Deep green: a quiet nod to NVIDIA's brand green, dark enough to read as text.
\definecolor{accent}{HTML}{3F6B00}  % THE accent: name, sections, links, bullets
\definecolor{ink}{HTML}{1A1A1A}     % body text (warm near-black, not pure 000)
\definecolor{soft}{HTML}{6E6E6E}    % secondary information (context, tech)
\definecolor{hair}{HTML}{C9C9C9}    % section rules
\definecolor{chipbg}{HTML}{EDF3E1}  % tag pill background
\definecolor{chipfg}{HTML}{3A5F0A}  % tag pill text

\hypersetup{colorlinks=true, urlcolor=accent, linkcolor=accent}
\AtBeginDocument{\color{ink}}

\pagestyle{empty}
\setlength{\parindent}{0pt}
\setlength{\parskip}{0pt}
\linespread{1.02}
\RaggedRight
\hyphenpenalty=10000
\exhyphenpenalty=10000
\emergencystretch=1em
\microtypesetup{expansion=false}

%% Section: letter-spaced small caps over a hairline
\titleformat{\section}
  {\large\scshape\lsstyle\color{accent}}{}{0pt}{}
  [\vspace{-5pt}{\color{hair}\rule{\textwidth}{0.4pt}}]
\titlespacing*{\section}{0pt}{13pt}{5pt}


%% Tag pill: \chip{CUDA}
\newcommand{\chip}[1]{%
  \tikz[baseline=(c.base)]{\node[inner xsep=4.5pt, inner ysep=2pt, rounded corners=2.5pt,
    fill=chipbg, text=chipfg, font=\footnotesize] (c) {#1};}\hspace{1pt}%
}


%% NB: at the head of a `p' cell, {\color{...}} shifts content down one line — always \textcolor.

%% ── 1. Education: stacked entries, no tags ───────────────────────────────
\newcommand{\edu}[4]{%
  \textbf{#1}\hfill #2\par\vspace{1.5pt}
  \textcolor{accent}{\textit{#3}}\par\vspace{1.5pt}
  {\small\color{soft}#4}\par
}

%% ── 2. Projects: prose (title first) / tags ──────────────────────────────────
%% Each row sizes its own tag column: measure the pills' natural width, cap it,
%% and let the prose take the rest.
\newlength{\projgapb}\setlength{\projgapb}{0.4cm}
\newlength{\projmaxtag}\setlength{\projmaxtag}{3.35cm}
\newlength{\projtagw}\newlength{\projprosew}
\newsavebox{\projtagbox}
%% #1 title (leads the prose, no final period), #2 prose, #3 pills.
\newcommand{\proj}[3]{%
  \sbox\projtagbox{\begin{varwidth}[t]{\projmaxtag}\raggedleft #3\end{varwidth}}%
  \setlength{\projtagw}{\wd\projtagbox}%
  \setlength{\projprosew}{\dimexpr\textwidth-\projgapb-\projtagw\relax}%
  \noindent
  \parbox[t]{\projprosew}{\RaggedRight \textbf{#1.} #2}%
  \hspace{\projgapb}%
  \usebox\projtagbox
  \par
}
\newcommand{\projgap}{\vspace{10pt}}
%% Clickable title that stays in ink: \titlelink{url}{Title}
\newcommand{\titlelink}[2]{\href{#1}{\textcolor{ink}{#2}}}

%% ── 3. Experience: role first, full-width prose ───────────────────────
%% #1 role, #2 date, #3 employer, #4 pills, #5 prose.
\newcommand{\job}[5]{%
  \textbf{#1}\hfill #2\par\vspace{1.5pt}
  \textcolor{accent}{#3}\hfill {#4}\par\vspace{4pt}
  #5\par
}

%% ── 4. Skills: label + pills ─────────────────────────────────────
\newlength{\skilllab}\setlength{\skilllab}{4.95cm}
\newcommand{\skillrow}[2]{%
  \noindent\parbox[t]{\skilllab}{\textbf{#1}}%
  \parbox[t]{\dimexpr\textwidth-\skilllab\relax}{\RaggedRight #2}\par\vspace{3.5pt}
}


\begin{document}

%% ── Header: full width, no photo ─────────────────────────────────────
{\color{accent}\fontsize{25}{30}\selectfont\scshape\textls[40]{Talvin Ackbaraly}}\\[6pt]
\textbf{Final-year internship — Developer Technology, CUDA \& GPU Computing · 6 months from February 2027}\\[5pt]
{\small
\href{mailto:talvin.ackbaraly@icloud.com}{talvin.ackbaraly@icloud.com} \enspace·\enspace
\href{tel:+33769389008}{+33 7 69 38 90 08} \enspace·\enspace
Paris, France\\[2pt]
\href{https://talvin-ackbaraly.com}{talvin-ackbaraly.com} \enspace·\enspace
\href{https://github.com/talvinckb}{github.com/talvinckb} \enspace·\enspace
\href{https://www.linkedin.com/in/talvin-ackbaraly}{linkedin.com/in/talvin-ackbaraly}}\\[7pt]
Final-year engineering student at EPITA (Image specialization), focused on GPU computing, real-time
graphics and performance. I enjoy the developer technology loop: profile an application, find what
limits it, rework it for the GPU, measure the gain, and explain it to other developers.

\section{GPU Computing \& Graphics}

\proj{\titlelink{https://www.talvin-ackbaraly.com/projects/irgpu/}{GPU motion detection}}{Real-time video background subtraction as a custom GStreamer plugin,
  with interchangeable C++ and CUDA backends. Profiled with Nsight to locate the bottlenecks, then
  optimized them: \textbf{$\times$24 speedup} over the CPU version.}
  {\chip{CUDA}\chip{C++}\chip{Nsight}}
\projgap
\proj{\titlelink{https://www.talvin-ackbaraly.com/projects/pogl/}{FluidSim: real-time 3D water simulation}}{SPH fluid physics running entirely on the GPU in
  GLSL compute shaders (density, pressure, viscosity), with GPU spatial hashing and a parallel
  bitonic sort for neighbor search. Screen-space water rendering with reflections and refraction.}
  {\chip{OpenGL}\chip{GLSL}\chip{Compute shaders}}
\projgap
\proj{Paper study: parallel tree of shapes}{In-depth study of \textit{A Tree of Shapes Computation
  Algorithm for Massively Parallel Architectures}: how a classically sequential image hierarchy is
  restructured for massively parallel hardware. Presented orally with my own slides.}
  {\chip{Parallel algorithms}\chip{GPU}}


\section{Experience}

\job{\titlelink{https://www.talvin-ackbaraly.com/projects/pfee-bnf/}{Final-Year Project — Computer Vision} \textnormal{\color{soft}(team of 4)}}{May 2026 -- ongoing}
    {EPITA $\times$ National Library of France (Gallica)}{\chip{Python}\chip{PyTorch}\chip{YOLO}}
    {Detection and classification of heritage illustrations in Gallica's digitized documents, via the
     IIIF API. In charge of detection: training and benchmarking models on a 10,000-image dataset
     annotated by the BnF, hyperparameter tuning, reaching \textbf{94.7\%~mAP50} on the validation set.}
\vspace{11pt}
\job{Full-Stack \& DevOps Developer \textnormal{\color{soft}(internship)}}{Sept. 2025 -- Jan. 2026}
    {Studiocanal}{\chip{Java}\chip{Spring Boot}\chip{AWS}\chip{GitLab CI}}
    {Built features for an internal rights, contracts and catalogue platform used by 100+ employees,
     within an Agile team. Set up the GitLab CI/CD pipeline for a new module and worked on its AWS
     infrastructure (EC2, RDS, IAM).}


\section{Other Projects}

\proj{\titlelink{https://www.talvin-ackbaraly.com/projects/alpr/}{License plate recognition}}{Classical image processing pipeline (filters, morphology, HOG) and
  Random Forest, \textbf{86\% precision}. Python vs C++ performance benchmark.}
  {\chip{C++}\chip{OpenCV}}
\projgap
\proj{\titlelink{https://github.com/talvinckb/3D-L-System}{3D L-system plant generator}}{Procedural plants from L-system grammars interpreted by a 3D
  turtle, drawn with instanced OpenGL rendering and custom GLSL shaders.}
  {\chip{OpenGL}\chip{GLSL}\chip{C++}}


\section{Education}

\edu{EPITA — Engineering School of Computer Science}{Sept. 2022 -- July 2027}
    {Master of Engineering, Image specialization — Paris}
    {Core curriculum: algorithms \& data structures, systems programming (C, Unix).\\
     Image specialization: GPU computing (CUDA), image processing \& synthesis, computer vision.}
\vspace{7pt}
\edu{Université du Québec à Chicoutimi (UQAC)}{Jan. -- May 2024}
    {Academic exchange semester — Québec, Canada}
    {}


\section{Skills}

\skillrow{GPU \& Graphics}{\chip{CUDA}\chip{Nsight}\chip{Profiling}\chip{OpenGL}\chip{GLSL}\chip{Compute shaders}}
\skillrow{Programming}{\chip{C++}\chip{C}\chip{Python}\chip{Java}\chip{SQL}}
\skillrow{Tools \& Libraries}{\chip{Linux}\chip{Git}\chip{CMake}\chip{GitLab CI/CD}\chip{Docker}\chip{GStreamer}\chip{Dear ImGui}\chip{OpenCV}\chip{PyTorch}}
\skillrow{Languages}{French (native), English (professional, TOEIC 845)}

\end{document}
